\section{验证：\texttt{int3} 自检}

\codefile{src/kernel/core/kmain.c}
\begin{lstlisting}[style=cstyle]
// kmain.c
idt_init();
printk("[init] idt ok\n");
__asm__ volatile("int3");   // 触发 #BP
printk("int3 returned, IDT works!\n");
\end{lstlisting}

预期输出：

\begin{tcblisting}{consolebox}
[init] idt ok
[isr] #BP Breakpoint vector=3 err=0x00000000 eip=0xC01xxxxx
int3 returned, IDT works!
\end{tcblisting}

如果看不到 ``int3 returned''，说明 \texttt{iret} 回不来——检查栈布局。

\subsection{GDB 逐指令验证}

\begin{tcblisting}{consolebox}
«\tps{(gdb)}» «\tin{b isr\_handler\_c}»
«\tps{(gdb)}» «\tin{c}»
Breakpoint, isr_handler_c (r=0xc01xxxxx)
«\tps{(gdb)}» «\tin{p r->int\_no}»
$1 = 3
«\tps{(gdb)}» «\tin{p/x r->eip}»
$2 = 0xc010abcd     # int3 指令的下一条地址（Trap 语义）
«\tps{(gdb)}» «\tin{p/x r->err\_code}»
$3 = 0x0             # int3 没有错误码
\end{tcblisting}

% ============================================================================

